\begin{tikzpicture}[x=1cm,y=1cm]
% функция F
\node[gv] (n1) at (0,0) {2};
\node[gv] (n2) at (-2,-1.6) {3};
\node[gv] (n5) at (2,-1.6) {7};
\node[gv] (n3) at (-3,-3.2) {4};
\node[gv] (n4) at (-1,-3.2) {6};
\node[gv] (n6) at (1,-3.2) {8};
\node[gv] (n7) at (3,-3.2) {10};
\node[gv] (n8) at (0,-4.8) {11};
% функция S
\node[gv] (n9) at (8,0) {13};
\node[gv] (n10) at (6.5,-1.6) {14};
\node[gv] (n11) at (9.5,-1.6) {16};
\node[gv] (n12) at (9.5,-3.2) {17};
\node[gv] (n13) at (9.5,-4.8) {18};
\node[gv] (n14) at (8,-6.4) {19};
% аргументы вызова F (дуги 16-2, 16-3, 16-7, 16-10)
\draw (n11.east) -- (10.6,-1.6) |- (-2,1.2);
\foreach \x/\n in {-2/n2,0/n1,2/n5,3/n7} {
    \node[junction] at (\x,1.2) {};
    \draw[flow] (\x,1.2) -- (\n.north);
}
% результат r и результат F
\draw[flow] (n3) -- (n8);
\draw[flow] (n4) -- (n8);
\draw[flow] (n6) -- (n8);
\draw[flow] (n14.west) -| (n8.south);
\draw[flow] (n8) -- (n13);
% аргументы вызова S
\draw[flow] (n7) to[out=40,in=180] (n9);
\draw[flow] (n7) -- (n11);
\draw[flow] (n7) -- (n12);
\draw[flow] (n12.east) -- (11.1,-3.2) |- (n9.east);
\draw[flow] (n12) -- (n11);
\draw[flow] (n12) to[out=-20,in=-70,looseness=6] (n12);
% результат t и результат S
\draw[flow] (n10) -- (n14);
\draw[flow] (n13) -- (n14);
\draw[flow] (n14) to[bend right=35] (n13);
\end{tikzpicture}
